\section{Introduction}

Video footage from match broadcasts and team cameras can be segmented into short clips of individual batting strokes, and automatically reading these clips to judge a player's technique and value is a practical and increasingly studied problem in coaching, scouting and broadcasting. The same stroke played by two different batters produces very different visual and statistical signatures, and turning raw footage into a single, trustworthy, numerically grounded analysis is where most existing systems fall short.

Most widely studied sports-analytics systems are built around a single narrow task for a single audience, such as forecasting how the players of a team will move \cite{ref1}, \cite{ref2}, spotting events in broadcast video \cite{ref3}, or training cooperative agents inside a simulated match \cite{ref4}, \cite{ref5}. A newer line of agentic systems, built from several coordinating agents, has shown promise beyond motion prediction, for example in supply-chain decision-making \cite{ref6}, \cite{ref7}, and within cricket itself to help fans assemble a fantasy team \cite{ref8}. However, each of these systems is built around one task and one answer format, provides no mechanism for checking its own numbers against a trusted source before reporting them, and offers no way to serve more than one kind of cricket stakeholder from the same underlying system. This matters because a report that sounds confident but was never checked against evidence can mislead the very people---coaches, scouts, medical staff---who rely on it to make a decision. To address this critical gap, we introduce CricketMind, a single agentic system that turns a batting-shot video, or a plain cricket question with no video at all, into a checked, stakeholder-specific analytics report.

The main contributions of this research are outlined as follows:

\begin{itemize}
    \item Development of CricketMind, an agentic architecture that accepts either a batting-shot clip or a text-only question. For a clip, a perception front end samples 15 frames, encodes them with a vision encoder and a temporal transformer, and classifies the shot, and a confidence gate with a bounded retry loop ensures that a low-confidence label is never passed downstream silently. A text-only question bypasses the perception path entirely.
    \item A shared, structured Shot DNA record and two sanctioned calculation engines, a deterministic statistics engine and a Monte Carlo what-if simulator, that hold exclusive authority to produce numbers, drawing evidence from two separate retrieval indexes: a precise ball-by-ball index and a broader historical index. Every other component must cite the engines' values rather than compute its own.
    \item A dynamic supervisor router that activates only the relevant subset of twelve stakeholder-facing agents, covering audiences from coaches and players to medical staff, scouts and broadcasters. A router, either keyword-based or a language model, proposes the agents, and fixed rules decide: agents whose evidence does not exist for the query are removed, and at most four are activated.
    \item A faithfulness gate, the Faith Check, that compares every figure an activated agent cites against the values the two engines produced, sends only the failing agents back to regenerate, and withholds any finding that still fails once a bounded regeneration budget is spent, before the verified findings are composed into the final, stakeholder-specific report.
    \item As an evaluation of the reference implementation, an execution profile whose cost follows a closed-form model with a worst case fixed when the graph is constructed, and a comparison of a keyword router with a locally hosted language-model router on paraphrased queries.
\end{itemize}
